%=============================================================================
%  TESSERA — Corporate Finance individual assignment
%  Gagandeep Singh · Masters in Artificial Intelligence with Business
%  SP Jain School of Global Management, Dubai
%
%  Compiles with pdfLaTeX on Overleaf. Upload the thirteen figures from the
%  screenshots/ folder alongside this file (or into a screenshots/ subfolder).
%=============================================================================
\documentclass[11pt,a4paper]{article}

\usepackage[T1]{fontenc}
\usepackage[utf8]{inputenc}
\usepackage[english]{babel}
\usepackage{lmodern}
\usepackage[margin=2.4cm,top=2.6cm,bottom=2.6cm]{geometry}
\usepackage{amsmath,amssymb}
\usepackage{graphicx}
\usepackage{booktabs}
\usepackage{array}
\usepackage{caption}
\usepackage{float}
\usepackage{fancyhdr}
\usepackage{enumitem}
\usepackage{eso-pic}
\usepackage{tikz}
\usetikzlibrary{calc}
\usepackage[hidelinks]{hyperref}

\graphicspath{{screenshots/}{figures/}{./}}

%--- Single-line black border on every page ----------------------------------
\AddToShipoutPictureBG{%
  \begin{tikzpicture}[remember picture,overlay]
    \draw[line width=0.5pt]
      ($(current page.north west)+(1.2cm,-1.2cm)$)
      rectangle
      ($(current page.south east)+(-1.2cm,1.2cm)$);
  \end{tikzpicture}%
}

%--- Running footer ----------------------------------------------------------
\pagestyle{fancy}
\fancyhf{}
\renewcommand{\headrulewidth}{0pt}
\fancyfoot[C]{\footnotesize Corporate Finance $\cdot$ Capital Budgeting and
  Investment Decision $\cdot$ TESSERA \\[2pt] \thepage}

\captionsetup{font=small,labelfont=bf,width=0.92\textwidth}
\setlength{\parskip}{0.45em}
\setlength{\parindent}{0pt}

\newcommand{\aed}[1]{AED~#1}

\begin{document}

%=============================================================================
%  TITLE PAGE
%=============================================================================
\begin{titlepage}
\thispagestyle{empty}
\centering
\vspace*{2.6cm}

{\scshape\large Corporate Finance}\\[0.5em]
\rule{0.62\textwidth}{0.4pt}\\[2.4cm]

{\Huge\bfseries TESSERA}\\[1.0em]

{\large Capital Budgeting Report}\\[0.6em]

{\itshape Should a studio own its intelligence, or keep renting it?}\\[2.2cm]

{\itshape TESSERA: An AI-Enabled Capital Budgeting and Investment
Decision Application}\\[2.6cm]

{\bfseries Student}\\[0.7em]
Gagandeep Singh\\[2.0cm]

{\bfseries Instructor:} [Instructor Name]

\vfill

{\footnotesize Submitted for the partial fulfilment of the requirements of the
Corporate Finance module of the Masters in Artificial Intelligence with
Business.}

\vspace*{0.8cm}
\end{titlepage}

%=============================================================================
%  CONTENTS
%=============================================================================
\tableofcontents
\newpage

%=============================================================================
\section{Company and Project Background}

Meridian Frame Studios is a 34-person animation and visual-effects studio in
Dubai Production City. Since 2024 its pipeline has depended on AI-assisted
work---frame interpolation, denoising, upscaling and generative
pre-visualisation---all running on rented GPU capacity billed by the hour. That
line has grown into one of the studio's largest operating expenses, and it rises
with every project won.

The studio is therefore asking whether to stop renting compute and own it
instead. The company is fictional, but every market parameter is calibrated to
publicly observable rates for accelerator hardware, colocation and regional cloud
pricing. The analysis is delivered as a working application, TESSERA, which
computes every figure in this report from one verified engine
(Figure~\ref{fig:hero}).

\begin{figure}[H]
  \centering
  \includegraphics[width=\textwidth]{fig01_hero_cold_boot.png}
  \caption{Landing view of the application. The cluster is bound to the model:
  27 of 32 modules are lit because peak utilisation is 85\%, and the headline
  figures are computed live rather than illustrated.}
  \label{fig:hero}
\end{figure}

\section{Description of the Investment Decision}

The decision is a mutually exclusive choice between four options over a
five-year horizon:

\begin{table}[H]
\centering
\small
\begin{tabular}{clrl}
\toprule
& \textbf{Alternative} & \textbf{Outlay (AED)} & \textbf{Life} \\
\midrule
A & Own --- full 32-GPU cluster        & 7,060,000 & 5 years \\
B & Rent --- 3-year reserved commitment &    45,000 & 3 years \\
C & Hybrid --- own 16, burst to cloud   & 3,830,000 & 5 years \\
--- & Do nothing --- continue on-demand & --- & --- \\
\bottomrule
\end{tabular}
\caption{The alternatives under appraisal.}
\end{table}

Alternative A is the proposal on the table. The remaining options exist because a
capital decision is only meaningful against the next-best use of the same money
(Figure~\ref{fig:decision}).

\begin{figure}[H]
  \centering
  \includegraphics[width=\textwidth]{fig03_decision_briefing.png}
  \caption{The decision and the four alternatives as presented in the
  application.}
  \label{fig:decision}
\end{figure}

\section{Project Assumptions and Data Sources}

Revenue is internal cost displacement plus external resale. Sixty per cent of
utilised hours substitute for on-demand rental at \aed{20.20}/GPU-hour in year~1;
the balance is resold regionally at \aed{14.70}. Both rates erode at 8\% a year,
the rate implied by the last two hardware generations. Capacity is
$32 \times 8{,}760$ hours at 96\% availability, utilised at 65\%, 80\%, 85\%,
84\% and 78\%.

The corporate tax rate is 9\%, the UAE federal rate. The discount rate is derived
rather than assumed: an unlevered beta of 1.25 is relevered by the Hamada
relation at the studio's 30\% debt weight,

\[
\beta_L = \beta_U\left[1 + (1-T)\frac{D}{E}\right] = 1.25\left[1 + 0.91 \times
\tfrac{0.30}{0.70}\right] = 1.7375,
\]

giving a CAPM cost of equity of $4.3\% + 1.7375 \times 5.6\% + 2.0\% = 16.03\%$
including a small-company premium, an after-tax cost of debt of 5.37\%, and

\[
\mathrm{WACC} = 0.70 \times 16.03\% + 0.30 \times 5.37\% = 12.83\%.
\]

All inputs are editable in the application (Figure~\ref{fig:assumptions}).

\begin{figure}[H]
  \centering
  \includegraphics[width=\textwidth]{fig04_assumption_console.png}
  \caption{The assumption console, including the derivation of the 12.83\%
  hurdle rate.}
  \label{fig:assumptions}
\end{figure}

\section{Relevant and Irrelevant Cash Flows}

Only incremental, after-tax cash flows enter the appraisal. Relevant flows are
the capital outlay, the displaced rental cost, external resale receipts, energy
and bandwidth, colocation, the incremental MLOps engineer, insurance, the
depreciation tax shield, salvage and working capital.

Excluded as irrelevant are existing payroll and premises, unchanged by the
decision; allocated head-office overhead, an accounting apportionment rather than
an incremental outflow; and financing cash flows, since
the cost of capital already carries the financing charge and deducting interest
again would double-count it.

\section{Treatment of Sunk Costs, Opportunity Costs, Working Capital,
Depreciation, Taxes and Salvage Value}

\textbf{Sunk cost.} A thermal and power feasibility study costing \aed{120{,}000}
was paid in Q1~2026. The cash is gone whichever option is chosen, so it is
excluded from all four alternatives.

\textbf{Opportunity cost.} The studio owns the 45\,m\textsuperscript{2} technical
room the cluster would occupy, so it carries no rent---but using it forfeits
\aed{86{,}000} a year of sublease income, charged as a fixed cost to
Alternatives A and C.

\textbf{Working capital.} \aed{500{,}000} is committed at $t_0$ against resale
receivables and spares, recovered in full at $t_5$, and not depreciated.

\textbf{Depreciation.} Equipment and installation of \aed{6{,}560{,}000} are
capitalised and written down straight-line over five years to a residual of
\aed{1{,}054{,}800}---\aed{1{,}101{,}040} a year. Depreciation is not a cash flow;
it enters only through the tax shield:
\[
\mathrm{OCF} = (R - C - D)(1 - T) + D.
\]

\textbf{Taxes and salvage.} Salvage is estimated at 18\% of original equipment
cost, \aed{1{,}054{,}800}, exactly equal to closing book value---so the base case
produces no disposal gain and no disposal tax. The application holds the
depreciation estimate separate from realised salvage, so changing either produces
a genuine taxable gain or loss rather than silently cancelling.

\begin{figure}[H]
  \centering
  \includegraphics[width=\textwidth]{fig05_cashflow_ledger.png}
  \caption{Initial, annual operating and terminal cash flows.}
  \label{fig:ledger}
\end{figure}

\section{Capital Budgeting Calculations and Results}

The initial outlay is $-(5{,}860{,}000 + 700{,}000 + 500{,}000) =
\aed{-7{,}060{,}000}$ (Figure~\ref{fig:ledger}). Discounting the resulting
series at 12.83\%:

\[
\mathrm{NPV} = \sum_{t=0}^{5} \frac{CF_t}{(1+r)^t} = \aed{-51{,}476},
\qquad
\mathrm{PI} = \frac{PV_{\text{inflows}}}{PV_{\text{outflows}}}
= \frac{7{,}008{,}524}{7{,}060{,}000} = 0.993.
\]

\begin{table}[H]
\centering
\small
\begin{tabular}{lr@{\hskip 2em}lr}
\toprule
\textbf{Measure} & \textbf{Value} & \textbf{Measure} & \textbf{Value} \\
\midrule
NPV @ 12.83\%        & \aed{$-$51{,}476} & Payback period      & 3.80 years \\
IRR                  & 12.54\%         & Discounted payback    & never \\
MIRR                 & 12.67\%         & ARR (avg.\ investment)& 15.65\% \\
Profitability Index  & 0.993           & ARR (initial)         & 8.44\% \\
\bottomrule
\end{tabular}
\caption{Decision metrics for Alternative A. Every measure fails, but narrowly.}
\end{table}

The NPV shortfall is 0.73\% of the outlay---inside the margin of error of the
assumptions themselves. A decision this close cannot be settled by the base case
(Figure~\ref{fig:metrics}).

\begin{figure}[H]
  \centering
  \includegraphics[width=\textwidth]{fig06_metrics_grid.png}
  \caption{Decision metrics with the runtime consistency checks that verify
  them.}
  \label{fig:metrics}
\end{figure}

Break-even analysis is decisive where the base case is not. Accounting
break-even rises from 45.9\% to 77.5\% utilisation across the five years and cash
break-even from 22.3\% to 43.9\%, but \emph{NPV} break-even requires 85.4\% peak
utilisation against a plan of 85.0\%. The project needs the cluster to run
essentially exactly as forecast, for five years, merely to be worth building
(Figure~\ref{fig:breakeven}).

\begin{figure}[H]
  \centering
  \includegraphics[width=\textwidth]{fig07_breakeven.png}
  \caption{Cash, accounting and NPV break-even thresholds.}
  \label{fig:breakeven}
\end{figure}

\section{Comparison of Investment Alternatives}

\begin{table}[H]
\centering
\small
\begin{tabular}{clrrrrr}
\toprule
& \textbf{Alt} & \textbf{NPV (AED)} & \textbf{EAA (AED)} & \textbf{IRR} &
\textbf{PI} & \textbf{Disc.\ payback} \\
\midrule
A & Own    & $-51{,}476$ & $-14{,}575$ & 12.54\%  & 0.993 & never \\
B & Rent   & 693{,}082   & 292{,}697   & 440.45\% & 16.402\textsuperscript{*} & 0.39 yr \\
C & Hybrid & \textbf{726{,}442} & 205{,}691 & 20.64\% & 1.190 & 4.14 yr \\
\bottomrule
\end{tabular}
\caption{Alternatives compared. \textsuperscript{*}Alternative B employs almost
no capital, so ratio measures are arithmetically large but economically
meaningless; the application suppresses them.}
\label{tab:compare}
\end{table}

Two complications underlie Table~\ref{tab:compare}. First,
\textbf{unequal lives}:
Alternative~B runs three years against five, so equivalent annual annuity,
\[
\mathrm{EAA} = \frac{\mathrm{NPV} \times r}{1-(1+r)^{-n}},
\]
is the correct adjustment---and it reverses the ranking, favouring B at
\aed{292{,}697} against C's \aed{205{,}691}. Second, EAA assumes indefinite
replication on identical terms, which for a locked price in a market falling 8\%
a year does enormous work. Held for five years, the same commitment yields an NPV
of \aed{553{,}267} and an EAA of \aed{156{,}656}---below Alternative~C.
Alternative~B's advantage is real but self-liquidating. For mutually exclusive
projects NPV governs, and \textbf{Alternative~C creates the most value}
(Figures~\ref{fig:hybrid} and~\ref{fig:suppressed}).

\begin{figure}[H]
  \centering
  \includegraphics[width=0.86\textwidth]{fig12_metrics_hybrid.png}
  \caption{Alternative C, the hybrid option, on the same measures.}
  \label{fig:hybrid}
\end{figure}

\begin{figure}[H]
  \centering
  \includegraphics[width=0.86\textwidth]{fig13_ratio_metrics_suppressed.png}
  \caption{Alternative B, with ratio metrics suppressed because the option
  employs almost no capital.}
  \label{fig:suppressed}
\end{figure}

\section{Sensitivity and Scenario Analysis}

Each driver moves across its own plausible range rather than by a uniform
percentage, because a blanket shock ranks drivers by the room they were given
rather than by how uncertain they genuinely are. Utilisation is the widest,
swinging NPV by \aed{2{,}985{,}394}. Its switching value is a utilisation factor
of 1.005---five-tenths of one per cent above plan (Figure~\ref{fig:tornado}).

\begin{figure}[H]
  \centering
  \includegraphics[width=\textwidth]{fig08_sensitivity_tornado.png}
  \caption{Tornado chart, switching values and the two-way grid.}
  \label{fig:tornado}
\end{figure}

All \textbf{nine of nine} drivers cross $\mathrm{NPV}=0$ inside their plausible
range---including the corporate tax rate, the least contentious assumption in the
model. That is not nine separate problems; it is one problem stated nine times.
Crossing the two widest drivers produces 121 combinations, of which 88 (73\%) are
loss-making. Scenario analysis gives a range from \aed{$-$4{,}073{,}264} to
\aed{2{,}416{,}615}---a spread of 0.92 times the outlay, wider than the
investment itself (Figure~\ref{fig:scenarios}).

\begin{figure}[H]
  \centering
  \includegraphics[width=\textwidth]{fig09_scenario_triptych.png}
  \caption{Best, base and worst case, with the resulting outcome range.}
  \label{fig:scenarios}
\end{figure}

\section{Financial and Non-Financial Risks}

The dominant financial risk is \textbf{price erosion}: the studio would buy a
depreciating asset into a market whose price is falling 8\% a year, and the
switching value is 7.73\%---the project fails if erosion is even slightly worse
than assumed. \textbf{Utilisation risk} follows: five idle GPUs each carry
\aed{205{,}000} of capital and \aed{42{,}469} of annual fixed cost against no
revenue at all (Figure~\ref{fig:idle}). \textbf{Technology obsolescence} is
acute, since a successor generation would compress both resale rates and salvage
simultaneously.

Non-financial risks include concentrating the entire platform on one incremental
engineer, the absence of any signed resale customer, and the data-residency case
for in-country capacity---a genuine strategic benefit the NPV does not capture.

\begin{figure}[H]
  \centering
  \includegraphics[width=\textwidth]{fig02_gpu_inspector_idle.png}
  \caption{An idle GPU inspected: full capital and fixed-cost share against zero
  revenue.}
  \label{fig:idle}
\end{figure}

\section{AI-Generated Insights and Evaluation of Those Insights}

The application embeds Claude (Opus~5) as an interpretive layer, deliberately
constrained: the response schemas contain \emph{no numeric fields at all}, so the
model receives verified figures as structured data and cannot introduce a number
of its own. Its recommendation is displayed beside a deterministic, rule-based
verdict computed in plain code from criteria fixed in advance, so the AI output
has something independent to be measured against
(Figures~\ref{fig:verdict} and~\ref{fig:comparison}).

\begin{figure}[H]
  \centering
  \includegraphics[width=\textwidth]{fig10_ai_verdict_vs_rules.png}
  \caption{Claude's recommendation beside the deterministic rule-based verdict.}
  \label{fig:verdict}
\end{figure}

\begin{figure}[H]
  \centering
  \includegraphics[width=\textwidth]{fig11_ai_comparison_explainer.png}
  \caption{Claude comparing the alternatives, and the on-demand metric
  explainer.}
  \label{fig:comparison}
\end{figure}

The two disagree, which is the useful part. The rules return \textsc{review
further} because Alternative~A's NPV falls inside their $\pm5\%$ marginal band;
Claude returns \textsc{reject}. I judge Claude correct and the rules
correct-but-irrelevant: the rules appraise Alternative~A in isolation, where
``too close to call'' is honest, while the real decision is a choice among
alternatives, and a better-scaled option for the same demand already exists.
Claude's identification of break-even rather than NPV as the decisive statistic
matches my own reading. I discount its confidence, however: it is confident about
a ranking resting on the same contestable utilisation forecast as every other
measure.

\section{Final Investment Recommendation}

\textbf{Reject Alternative A; adopt Alternative C.} Alternative~A fails every
capital test, and its NPV break-even of 85.4\% against a plan of 85.0\% leaves no
operational headroom whatsoever. Alternative~C serves the same underlying demand
for \aed{3{,}830{,}000}, returns an NPV of \aed{726{,}442}, a PI of 1.190 and an
IRR of 20.64\%, and works for the opposite reason: it buys only capacity that is
structurally saturated and pays status-quo prices for the peak, where owning
confers no advantage.

Three conditions attach. Size the owned tranche to contracted or structurally
always-on demand rather than to forecast peak; secure committed offtake from at
least two regional clients before releasing capital for the resale-dependent
portion; and set a review trigger so that observed erosion above 10\% a year
forces re-appraisal before any second tranche.

\section{Limitations of the Analysis}

The model is deterministic. It reports a range but attaches no probabilities to
it, and a Monte Carlo treatment would express the utilisation risk far better
than three discrete cases. Correlations between drivers are imposed by
judgement rather than estimated from data. Residual value beyond
year five is set to salvage with no continuing value recognised, understating
Alternative~A if the hardware remains serviceable. The resale market is assumed
to exist at the rates modelled, and no customer has validated it. Finally, real options---the ability to defer, to stage the purchase
or to abandon---carry genuine value that a static NPV cannot capture, and staging
is precisely what Alternative~C represents.

\section*{Verification}
\addcontentsline{toc}{section}{Verification}

Every figure was checked three ways: 69 unit tests whose expected values were
derived by hand; an independent re-implementation of the whole model in Python
against \texttt{numpy-financial}, agreeing on all 79 compared quantities; and 28
runtime consistency checks. One genuine disagreement was found and fixed---the
original solver returned no IRR for a non-conventional series because its
bracketing had failed rather than because it had detected multiple roots. The
solver now sweeps for every root; that series has two ($-37.79\%$ and
$+439.01\%$), and the application reports both.

\section*{Declaration of AI Tools Used}
\addcontentsline{toc}{section}{Declaration of AI Tools Used}

This project was built end to end using \textbf{Claude} (Anthropic) via Claude
Code, which was used to write the application, the financial engine, the test
suite and the independent Python verification, and to draft this report. Claude
Opus~5 is additionally embedded in the delivered application as its advisory
layer. No other large language model was used. All financial figures are computed
by deterministic code and independently verified; the AI layer interprets those
figures and cannot generate them. The final judgement, the choice of scenario, the
assumption set and the recommendation are my own.

\section*{References}
\addcontentsline{toc}{section}{References}

\begin{enumerate}[leftmargin=1.4em,itemsep=0.2em]
  \item Brealey, R.A., Myers, S.C. and Allen, F. (2023) \textit{Principles of
    Corporate Finance}. 14th edn. New York: McGraw-Hill.
  \item Ross, S.A., Westerfield, R.W. and Jaffe, J. (2022) \textit{Corporate
    Finance}. 13th edn. New York: McGraw-Hill.
  \item Damodaran, A. (2026) \textit{Betas by Sector} and \textit{Cost of Capital
    by Industry}. NYU Stern School of Business. Available at:
    \url{https://pages.stern.nyu.edu/~adamodar/}.
  \item United Arab Emirates (2022) \textit{Federal Decree-Law No.\ 47 of 2022 on
    the Taxation of Corporations and Businesses}. Ministry of Finance, UAE.
  \item Hamada, R.S. (1972) `The effect of the firm's capital structure on the
    systematic risk of common stocks', \textit{The Journal of Finance}, 27(2),
    pp.\ 435--452.
  \item Anthropic (2026) \textit{Claude Developer Documentation}. Available at:
    \url{https://docs.anthropic.com/}.
  \item Application source code and verification suite:
    \url{https://github.com/gagu00000/tessera-capital-budgeting}.
  \item Live application: \url{https://tessera-vert-nine.vercel.app}.
\end{enumerate}

\end{document}
